\documentclass[11pt,a4paper]{article}

\usepackage[utf8]{inputenc}
\usepackage[T1]{fontenc}
\usepackage{amsmath,amssymb,bm}
\usepackage{booktabs}
\usepackage[margin=2.5cm]{geometry}
\usepackage{hyperref}
\usepackage{enumitem}
\usepackage{listings}
\usepackage{xcolor}

\hypersetup{
    colorlinks=true,
    linkcolor=blue,
    citecolor=blue,
    urlcolor=blue
}

\lstset{
    basicstyle=\ttfamily\footnotesize,
    columns=fullflexible,
    breaklines=true,
    breakatwhitespace=true,
    frame=single,
    framerule=0.3pt,
    rulecolor=\color{gray!50},
    backgroundcolor=\color{gray!5},
    xleftmargin=0.5em,
    xrightmargin=0.5em,
    aboveskip=0.6em,
    belowskip=0.6em
}

\newcommand{\mb}[1]{\mathbf{#1}}
\newcommand{\Heff}{\mb{H}_{\text{eff}}}
\newcommand{\hth}{\mb{h}_{\text{th}}}
\newcommand{\Ms}{M_s}
\newcommand{\Aex}{A_{\text{ex}}}
\newcommand{\Keff}{K_{\text{eff}}}
\newcommand{\Vcell}{V_{\text{cell}}}
\newcommand{\tCo}{t_{\text{Co}}}
\newcommand{\Rth}{R_{\text{th}}}
\newcommand{\jcur}{j}
\newcommand{\kB}{k_B}

\title{Stochastic LLGS Solver for Thermal Skyrmion Dynamics:\\
Theory, Implementation Plan, and Reasoning}
\author{Rui Barreira}
\date{May 2026}

\begin{document}
\maketitle

\tableofcontents
\bigskip\hrule\bigskip

%======================================================================
\section{Motivation and questions to answer}
%======================================================================

The deterministic LLGS simulator in \texttt{src/simulator/} reproduces
the $T = 0$ dynamics of a SAF skyrmion under SOT drive but cannot
address the following physics questions:

\begin{enumerate}[leftmargin=2em]
    \item[\textbf{Q1}] At room temperature, can the same drive current
    sustain the same skyrmion velocity as at $T = 0$, or must the
    drive be reduced?
    \item[\textbf{Q2}] At high drive velocity, is the directional
    control preserved? In a SAF the topological Hall force cancels
    between layers; thermal fluctuations partially break this
    cancellation. How does the skyrmion Hall angle $\theta_H$ and the
    transverse position variance $\sigma_y(t)$ behave as $T$ and
    $\jcur$ grow?
    \item[\textbf{Q3}] Does the skyrmion size affect velocity and
    inter-skyrmion interactions? The skyrmion radius $R$ is set by
    the balance of $\Aex$, $D$, $\Keff$ and (in our parameterization)
    the applied field $H_z$. The radius enters the pair potential
    through the domain-wall width $\Delta = \sqrt{\Aex/\Keff}$.
    \item[\textbf{Q4}] Is there an optimal temperature that maximizes
    drift speed while preserving the skyrmion against thermal
    annihilation?
\end{enumerate}

Answering these requires solving the \emph{stochastic} LLGS equation
of Brown~\cite{brown1963} on the SAF lattice, with Joule heating
coupling the drive current density to the local temperature.

%======================================================================
\section{Theory}
%======================================================================

\subsection{Brown's stochastic LLG equation}

Magnetic fluctuations at finite temperature are modeled by augmenting
the deterministic effective field with a Gaussian white-noise field
$\hth(\mb{r}, t)$ that satisfies the fluctuation--dissipation theorem:
%
\begin{equation}\label{eq:sllg}
\frac{d\mb{m}}{dt} = -\gamma' \, \mb{m} \times (\Heff + \hth)
- \gamma' \alpha \, \mb{m} \times \big[\mb{m} \times (\Heff + \hth)\big]
+ \bm{\tau}_{\text{SOT}},
\end{equation}
%
with $\gamma' = \gamma / (1 + \alpha^2)$. The thermal field has zero
mean and is delta-correlated in space and time:
%
\begin{equation}\label{eq:fdt}
\langle h_{\text{th},i}^{(\ell)}(\mb{r}, t)\, h_{\text{th},j}^{(\ell')}(\mb{r}', t') \rangle
= \sigma^2 \, \delta_{ij}\, \delta_{\ell\ell'}\,
\delta(\mb{r} - \mb{r}')\, \delta(t - t'),
\end{equation}
%
with $\ell \in \{\text{top}, \text{bot}\}$ indexing the two SAF
layers. The noise variance is fixed by the FDT in the form
appropriate to the simulator's units (effective field in Tesla,
$\gamma$ in rad/s/T):
%
\begin{equation}\label{eq:sigma}
\boxed{\;
\sigma^2 = \frac{2\, \alpha\, \kB T}{\gamma\, \Ms\, \Vcell},
\qquad
\Vcell = a^2 \, \tCo.
\;}
\end{equation}
%
A short derivation of~\eqref{eq:sigma} is given in
Appendix~\ref{app:fdt}. The two layers of the SAF have independent
thermal baths.

\subsection{Stratonovich interpretation}

The noise in~\eqref{eq:sllg} multiplies $\mb{m}$ through the cross
product, so this is a multiplicative-noise SDE. The constraint
$|\mb{m}|=1$ is preserved only under the \emph{Stratonovich}
interpretation; an It\^{o} interpretation introduces a spurious drift
that violates the norm. We adopt Stratonovich throughout. The
practical consequence is that any midpoint or predictor-corrector
scheme that evaluates the noise term at both the predictor and
corrector stages \emph{with the same noise sample} converges to the
Stratonovich solution, and renormalization at each step keeps the
constraint to machine precision.

\subsection{Drive: SOT plus Joule heating}

The drive is the existing Slonczewski damping-like plus field-like
spin--orbit torque already encoded in \texttt{src/simulator/integrator.py}:
%
\begin{equation}
\bm{\tau}_{\text{SOT}} =
\gamma'\, \big[ (H_{DL} + \alpha H_{FL})\, \mb{m} \times \mb{s}
+ (H_{FL} - \alpha H_{DL})\, \mb{s} \big],
\end{equation}
%
with $\mb{s} = (\hat{\mb{p}} \times \mb{m})/|\hat{\mb{p}} \times \mb{m}|$,
$H_{DL} = c_{DL}\, \jcur$, $H_{FL} = c_{FL}\, \jcur$.

Joule self-heating raises the lattice temperature above the substrate
temperature:
%
\begin{equation}\label{eq:joule}
T(\jcur) = T_{\text{sub}} + \Rth\, \jcur^2,
\end{equation}
%
where $\Rth$ is a stack-dependent thermal resistance. We carry $\Rth$
as a free, user-specified parameter. No literature-anchored value for
this Pt/Co/Ru SAF is assumed; sensitivity to $\Rth$ is part of the
final output. Through~\eqref{eq:joule}, increasing $\jcur$ both
strengthens the drive and raises $T$, so the trade-off in Q4 is
non-trivial.

\subsection{Observables}

\begin{itemize}[leftmargin=2em]
    \item \textbf{Drift velocity}
    $\mb{v} = \langle d\mb{r}_c / dt \rangle$ from
    \texttt{skyrmion\_center} unwrapped against periodic boundaries.
    \item \textbf{Hall angle}
    $\theta_H = \arctan(\langle v_y \rangle / \langle v_x \rangle)$
    measured on the second half of the trajectory.
    \item \textbf{Transverse variance}
    $\sigma_y^2(t) = \langle (y_c(t) - \langle y_c(t)\rangle)^2 \rangle$.
    \item \textbf{Survival probability}
    $P_{\text{surv}}(T, \jcur) = \Pr(|Q(t_{\max})| > 0.5)$ across the
    ensemble.
    \item \textbf{Pair potential} $V(r)$ from short-time forced-pair
    runs at $\jcur = 0$, $T = 300$~K (Q3).
\end{itemize}

%======================================================================
\section{Numerical scheme}
%======================================================================

\subsection{Stochastic Heun predictor--corrector}

For Stratonovich SDEs with multiplicative noise, the stochastic Heun
scheme is the standard order-1.0 method. For the SAF pair
$(\mb{m}_t, \mb{m}_b)$ at step $n$, per layer:
%
\begin{enumerate}[leftmargin=2em]
    \item Sample a noise increment $\hth$ for each lattice cell with
    cell-wise standard deviation $\sigma / \sqrt{\Delta t}$. The
    \emph{same} sample is reused in both stages below.
    \item Compute $\Heff$ from the existing
    \texttt{effective\_field} (no noise inside the field assembly).
    Form $H_{\text{total}} = \Heff + \hth$ and call
    $f \equiv \texttt{llgs\_rhs}(\mb{m}, H_{\text{total}}, p)$.
    The $\mb{m}\times(\cdot)$ and $\mb{m}\times(\mb{m}\times(\cdot))$
    structure inside \texttt{llgs\_rhs} distributes the noise across
    the precession and damping terms, matching PysLLG's \texttt{Gfunc}
    form $\propto -\mb{m}\times(\hth + \alpha\, \mb{m}\times\hth)$.
    \item Predictor: $\tilde{\mb{m}} = \mb{m} + \Delta t \, f$.
    \emph{Do not} renormalize the predictor.
    \item Corrector (same noise sample, $\Heff$ recomputed at
    $\tilde{\mb{m}}$):
    $\mb{m}^{n+1} = \mb{m} + \tfrac{\Delta t}{2}
    \big[ f(\mb{m}, \Heff(\mb{m}) + \hth) +
          f(\tilde{\mb{m}}, \Heff(\tilde{\mb{m}}) + \hth) \big]$.
    \item Before normalizing, log
    $\max_{i,j} \big|\,\|\mb{m}^{n+1}_{ij}\| - 1\,\big|$ and raise
    \texttt{RuntimeError} if it exceeds a user-supplied tolerance
    \texttt{tol\_norm}. Then renormalize once.
\end{enumerate}
%
Single end-of-step renormalization matches PysLLG and the standard
Stratonovich-Heun derivation; renormalizing the predictor introduces
an $O(\sigma^2 \Delta t)$ bias to the corrector evaluation that the
single-end version avoids. The two SAF layers carry independent
noise draws; within a step, each layer's draw is shared by its own
predictor and corrector.

\subsection{Time step}

We retain the deterministic $\Delta t = 5\times 10^{-14}$~s as a
starting point and verify a posteriori that the equilibrium
$\langle m_z^2 \rangle$ matches the Langevin prediction to within
5\%. If it does not, halve $\Delta t$ and re-run the Langevin
validation. No automatic step-size adaptation is used; the time step
is user-supplied.

\subsection{Reproducibility}

Each trajectory consumes one explicit \texttt{numpy.random.Generator}
constructed from a user-supplied integer seed. Ensembles fold the
trajectory index into the seed (e.g., \texttt{seed\_base + traj\_idx})
so that any single trajectory is byte-reproducible without rerunning
the rest of the ensemble.

%======================================================================
\section{Implementation plan}
%======================================================================

\subsection{Scope and constraints}

\begin{itemize}[leftmargin=2em]
    \item All new code lives in \texttt{src/stochastic\_llgs/}.
    Existing files are \emph{read} and imported but never modified.
    \item Reuse \texttt{effective\_field} from
    \texttt{src/simulator/fields.py}, \texttt{llgs\_rhs} from
    \texttt{src/simulator/integrator.py}, and
    \texttt{effective\_field\_demag\_pair} from
    \texttt{src/phase\_diagram/fields\_demag.py}. Reuse
    \texttt{skyrmion\_center} and \texttt{skyrmion\_diameter} from
    \texttt{src/simulator/analysis.py}, \texttt{topological\_charge}
    from \texttt{src/simulator/main.py}, and \texttt{saf\_skyrmion}
    from \texttt{src/simulator/initial\_conditions.py}.
    \item No \texttt{argparse}. Every runnable module has a
    \texttt{main()} that opens with a labeled User Configuration
    block of snake\_case Python variables.
    \item No default argument values anywhere. Every function
    parameter is provided explicitly by the caller. Missing or
    invalid inputs raise \texttt{RuntimeError} with a precise message
    naming the offending field.
    \item No silent failure paths. No \texttt{try/except} that logs
    and continues. No \texttt{dict.get(key, fallback)}. The
    end-of-step norm-drift check logs
    $\max\big|\,\|\mb{m}\| - 1\,\big|$ and raises
    \texttt{RuntimeError} if it exceeds a user-supplied
    \texttt{tol\_norm}, instead of letting renormalization mask a
    too-coarse $\Delta t$.
    \item Units SI/Tesla throughout.
    \item Backend strategy is \emph{hybrid}: NumPy stochastic Heun
    first (Phase 1--5). A JAX/\texttt{diffrax} port is implemented
    only if Phase 5 wall time exceeds the working budget (Phase 6,
    optional).
\end{itemize}

\subsection{File layout (new code only)}

\begin{lstlisting}
src/stochastic_llgs/
  __init__.py
  parameters_thermal.py     # attach T, R_th, sigma to a p that
                            # the caller has already populated
  thermal_field.py          # sample_thermal_field(rng, shape,
                            #                     sigma, dt)
  integrator_sllg.py        # heun_stochastic_step for the SAF
                            # pair, with optional demag kernels
  joule_heating.py          # T_of_j(j, T_sub, R_th)
  diagnostics.py            # PBC-unwrapped trajectory, Q track,
                            # annihilation flag, Hall angle fit
  ensemble.py               # ProcessPoolExecutor ensemble driver
  io.py                     # NPZ writers
  validation/
    __init__.py
    macrospin.py
    test_langevin.py
    test_brown_reversal.py
    test_equipartition.py
    test_t0_limit.py
  production/
    __init__.py
    run_single.py
    scan_tj.py
    scan_radius.py
    pair_potential.py
  jax_backend/              # Phase 6 only; absent until needed
    fields_jax.py
    integrator_diffrax.py
\end{lstlisting}

\subsection{Phase 1 --- Noise generator and macrospin solver}

\begin{enumerate}[leftmargin=2em]
    \item \texttt{parameters\_thermal.attach\_thermal(p, T, R\_th, seed)}
    attaches \texttt{p.T, p.R\_th, p.seed, p.V\_cell, p.k\_B,
    p.sigma\_noise}. All four arguments required; validate finite and
    positive; raise \texttt{RuntimeError} naming any offender.
    \item \texttt{thermal\_field.sample\_thermal\_field(rng, shape,
    sigma, dt)} returns Gaussian noise of shape
    \texttt{(ny, nx, 3)} with std $\sigma/\sqrt{\Delta t}$. The
    \texttt{rng} is supplied by the caller; constructing one
    silently is forbidden.
    \item \texttt{integrator\_sllg.heun\_stochastic\_step(m\_top,
    m\_bot, dt, p, rng, demag\_kernels)} implements the predictor and
    corrector with the same noise sample for both layers and both
    stages.
\end{enumerate}

Phase 1 gates on the macrospin Langevin check
(Section~\ref{sec:val}).

\subsection{Phase 2 --- Lattice noise and equipartition}

Phase 2 verifies that the noise field has the correct cell-resolved
amplitude. A $16 \times 16$ ferromagnetic patch at $T = 10$~K is
equilibrated, then the mode-resolved energy of low-$k$ magnons is
compared to the Rayleigh--Jeans equipartition $\langle E_k \rangle =
\kB T / 2$. Passing within 20\% on modes with $k < \pi/(4a)$ is the
gate.

\subsection{Phase 3 --- SAF, $T=0$ reproducibility (split)}

The stochastic stepper accepts an optional
\texttt{demag\_kernels} argument. When \texttt{None}, the stepper
calls \texttt{effective\_field} (no demag) for each layer. When a
populated dict from \texttt{precompute\_demag\_kernels(p)} is
supplied, the stepper calls \texttt{effective\_field\_demag\_pair}
from \texttt{src/phase\_diagram/fields\_demag.py}.

\textbf{Phase 3a (no-demag baseline).} Set $T=0$,
\texttt{demag\_kernels=None}, and run the same SOT-driven analysis
as \texttt{src/simulator/analysis.py:run\_analysis}. \textbf{Pass:}
equilibrium diameter, drift velocity, Hall angle, $Q$ each agree
with the deterministic RK4 to within 2\%.

\textbf{Phase 3b (demag pathway).} Set $T=0$ and pass
\texttt{precompute\_demag\_kernels(p)}. The reference is generated
on the fly by driving the existing \texttt{rk4\_step} with
\texttt{effective\_field\_demag\_pair} for the same duration; no
existing code is modified. \textbf{Pass:} 2\% on each observable.

\subsection{Phase 4 --- Diagnostics and Joule heating}

\begin{itemize}[leftmargin=2em]
    \item \texttt{diagnostics.track\_skyrmion} unwraps the center
    against periodic boundaries via the phase of the circular mean
    $\langle e^{i 2\pi x_c / L}\rangle$ before differencing.
    \item \texttt{diagnostics.detect\_annihilation(Q\_history,
    q\_threshold, k\_consecutive)} requires both
    \texttt{q\_threshold} and \texttt{k\_consecutive} from the
    caller. It returns the first step index $n$ such that
    $|Q[m]| < $\,\texttt{q\_threshold} holds for all $m \in
    [n - k_c + 1, n]$, or $-1$ if the skyrmion survives. The
    consecutive-step guard suppresses false positives from
    finite-$T$ topological-charge fluctuations.
    \item \texttt{joule\_heating.T\_of\_j(j, T\_sub, R\_th)}
    implements~\eqref{eq:joule} as a uniform global temperature
    (one scalar $T$ per run); spatially resolved heating is not
    modeled in this phase.
\end{itemize}

\subsection{Phase 5 --- Production runs (NumPy)}

Shared configuration: lattice $256 \times 256$,
$N_{\text{ens}} = 10$ per cell, $\Delta t = 5\times 10^{-14}$~s,
$t_{\max} = 2$~ns. Joule heating is uniform and global,
$T = T_{\text{sub}} + R_{\text{th}} \jcur^2$.

\begin{itemize}[leftmargin=2em]
    \item \texttt{run\_single.py}: one $(T_{\text{sub}}, \jcur)$
    trajectory. Demag is a user-set boolean in the User
    Configuration block.

    \item \texttt{scan\_tj.py} (Q1, Q2, Q4): \textbf{demag OFF}.
    Grid $T_{\text{sub}} \in \{200, 250, 300, 350, 400\}$~K $\times$
    $\jcur \in \{1, 2, 4, 8, 16\} \times 10^{11}$~A/m$^2$. Per cell:
    $\langle v\rangle, \sigma_v, \theta_H, \sigma_{\theta_H}$, and
    $P_{\text{surv}} = (\#\text{alive at }t_{\max})/N_{\text{ens}}$
    using the \texttt{detect\_annihilation} guard.

    \item \texttt{scan\_arrhenius.py} (stability, dedicated): demag
    OFF, $\jcur = 0$, $N_{\text{ens}} = 30$. Sweep $T \in \{300,
    350, 400, 450, 500\}$~K (high $T$ chosen so annihilation events
    occur within $t_{\max}$). Each trajectory runs to $t_{\max}$ or
    annihilation, whichever comes first. Fit $\log \tau$ vs $1/T$
    and report $\Delta E_{\text{barrier}}$ and the Brown prefactor
    $\tau_0$. This is the canonical SAF skyrmion thermal-stability
    figure.

    \item \texttt{scan\_radius.py} (Q3): \textbf{demag ON}. At
    $T = 300$~K, sweep $(D, H_z)$ to span
    $R \in \{40, 60, 80, 120\}$~nm. Report $v(R), \theta_H(R),
    P_{\text{surv}}(R)$.

    \item \texttt{pair\_potential.py} (Q3): \textbf{demag ON}.
    Forced-pair runs at $T = 300$~K, $\jcur = 0$. Reconstruct
    $V(r)$ from short-time radial force measurements.
\end{itemize}

\paragraph{Why demag-off for the stability scans, demag-on for
radius.} The thin-film shape-anisotropy correction $-\mu_0 \Ms$ is
already folded into the precomputed $C_{\text{anis}}$ in
\texttt{parameters.py:132}, capturing the leading-order
uniform-magnetization demag. The non-uniform residual from
in-plane spin divergence in the domain wall shifts the skyrmion
radius by 5--10\%, $\langle v \rangle$ and $\theta_H$ by $<5\%$,
and does not change the qualitative trend of $P_{\text{surv}}(T,
\jcur)$ or $\tau(T)$. For Q3 the radius itself is the observable,
so the demag pathway is enabled there.

\subsection{Phase 6 --- JAX/diffrax port (only if needed)}

Triggered only if the Phase 5 wall time on the production grid
exceeds the user's working budget. Port \texttt{effective\_field}
and the demag convolution to JAX; use \texttt{diffrax.MultiTerm}
with the \texttt{GeneralShARK} solver and a
\texttt{VirtualBrownianTree} keyed by seed; \texttt{jax.vmap} over
the ensemble axis. Regress against a single NumPy cell before
swapping.

%======================================================================
\section{Validation gates}\label{sec:val}
%======================================================================

Every gate must produce its named NPZ artifact and pass its
quantitative criterion before the next phase begins.

\subsection{Langevin macrospin (Phase 1)}

Single spin in a Zeeman field $B \hat{z}$, no anisotropy, $\alpha=1$.
For each $T \in \{50, 100, 200, 300, 500\}$~K and several $B$ values,
run $\geq 200$ trajectories of length
$\gtrsim 100 / (\alpha \gamma B)$. Compare the time- and
ensemble-averaged $\langle m_z \rangle$ to the Langevin function
$L(\mu B / \kB T)$ with $\mu = \Ms \Vcell$.
\textbf{Pass:} RMS$(\langle m_z\rangle_{\text{sim}} - L)/L < 0.05$.

\subsection{Brown reversal time (Phase 1)}

Uniaxial macrospin with $K = K_{\text{top}}$, no field. For
temperatures yielding $K\Vcell/\kB T \in \{3, 5, 8\}$ (higher
barriers are uncomputably slow), measure the mean first-passage time
across $\geq 100$ trajectories and compare to Brown's high-barrier
formula
%
\begin{equation}
\tau \;\sim\; \frac{(1+\alpha^2)}{\alpha\gamma}\,
\sqrt{\frac{\pi \kB T}{K \Vcell}}\,
\exp\!\Big(\frac{K \Vcell}{\kB T}\Big).
\end{equation}
%
\textbf{Pass:} $|\log_{10}(\tau_{\text{sim}} / \tau_{\text{Brown}})|
< 0.5$ at each barrier.

\subsection{Spin-wave equipartition (Phase 2)}

$16 \times 16$ ferromagnetic patch, single layer (RKKY $= 0$), no
DMI, $K = 0$, weak $H_z$. Equilibrate at $T = 10$~K, then accumulate
$\langle |\delta m_k|^2 \omega_k^2 \rangle$. \textbf{Pass:} match to
$\kB T / 2$ within 20\% for modes with $k < \pi/(4a)$.

\subsection{$T=0$ deterministic limit (Phase 3a, no demag)}

Set $T = 0$, \texttt{demag\_kernels = None}, and run the same
SOT-driven analysis as \texttt{src/simulator/analysis.py:run\_analysis}.
\textbf{Pass:} equilibrium diameter, drift velocity, Hall angle, and
$Q$ each agree with the deterministic RK4 to within 2\%.

\subsection{$T=0$ deterministic limit (Phase 3b, with demag)}

Same setup but with \texttt{precompute\_demag\_kernels(p)} passed
in. The reference is generated by driving the existing
\texttt{rk4\_step} with \texttt{effective\_field\_demag\_pair} for
the same duration; no edits to existing code. \textbf{Pass:} 2\% on
each observable.

%======================================================================
\section{Run-script conventions (no argparse)}
%======================================================================

Every runnable module's \texttt{main()} opens with a labeled User
Configuration block of snake\_case Python variables. The user edits
the file to reconfigure a run. SLURM array indices, when needed, are
read from environment variables with explicit \texttt{os.environ[key]}
so a missing variable raises \texttt{KeyError} rather than silently
defaulting.

\begin{lstlisting}[language=python]
def main():
    # ================ User Configuration ================
    t_kelvin    = 300.0
    j_current   = 4.0e11        # A/m^2
    r_th        = 1.0e-15       # K m^4 / A^2
    t_max       = 2.0e-9        # s
    dt          = 5.0e-14       # s
    seed_base   = 17
    n_ens       = 10
    n_workers   = 8
    nx, ny      = 256, 256
    # ============ End User Configuration =================
    ...
\end{lstlisting}

A sibling \texttt{.sh} SLURM wrapper loads modules and invokes
\texttt{python -m src.stochastic\_llgs.production.scan\_tj}; per-run
knobs (grids, ensemble size, lattice size) live inside the Python
file itself.

%======================================================================
\section{Verification plan (end-to-end)}
%======================================================================

\begin{enumerate}[leftmargin=2em]
    \item \textbf{Langevin gate.} Run
    \texttt{validation/test\_langevin.py}; confirm RMS criterion.
    \item \textbf{Brown gate.}
    \texttt{validation/test\_brown\_reversal.py};
    confirm log-time criterion.
    \item \textbf{Equipartition gate.}
    \texttt{validation/test\_equipartition.py};
    confirm low-$k$ modes within 20\%.
    \item \textbf{$T=0$ gates.}
    \texttt{validation/test\_t0\_limit\_nodemag.py} and
    \texttt{validation/test\_t0\_limit\_demag.py}; confirm 2\%
    agreement against their respective deterministic references.
    \item \textbf{Single-cell production.}
    \texttt{production/run\_single.py} at $(T_{\text{sub}}, \jcur) =
    (300~\text{K}, 4\times 10^{11}~\text{A/m}^2)$. Inspect the
    trajectory; confirm drift, jitter, and survival.
    \item \textbf{Cost calibration.} From the single-cell wall time,
    estimate total cost of the production grid. If above budget,
    trigger Phase 6.
    \item \textbf{Production scans.}
    \texttt{production/scan\_tj.py},
    \texttt{production/scan\_arrhenius.py},
    \texttt{production/scan\_radius.py},
    \texttt{production/pair\_potential.py}.
    \item \textbf{Answers.} A generated notebook
    \texttt{output/stochastic\_llgs/answers.md} reads the NPZ files
    and answers Q1--Q4 in the same order as the introduction.
    Q4 is answered jointly from \texttt{scan\_tj} (driven survival)
    and \texttt{scan\_arrhenius} (zero-drive barrier).
\end{enumerate}

%======================================================================
\section{Risks and open items}
%======================================================================

\begin{itemize}[leftmargin=2em]
    \item $\Rth$ has no literature-anchored value for this Pt/Co/Ru
    SAF stack. We carry it as a user-supplied parameter and report
    sensitivity to it in the final scan.
    \item $N_{\text{ens}} = 10$ may be insufficient for cells with
    $P_{\text{surv}}$ near 0.5; the script raises a warning and we
    increase $N_{\text{ens}}$ to 30 for tail cells.
    \item $\Delta t = 5\times 10^{-14}$~s is borrowed from the
    deterministic simulator; the Langevin gate provides the
    quantitative check that this $\Delta t$ is small enough at high
    temperature. If it fails at $T = 500$~K, halve $\Delta t$ for the
    full pipeline.
    \item Production cost on $256 \times 256$ with demag, $N_{\text{ens}}
    = 10$, $25$ $(T, \jcur)$ cells dominates the budget; the
    single-cell calibration in step 5 of the verification plan
    decides whether to trigger Phase 6.
\end{itemize}

%======================================================================
\appendix
\section{Derivation of the thermal noise variance}\label{app:fdt}
%======================================================================

Starting from the Landau--Lifshitz equation in SI form with
$\Heff$ in Tesla and $\gamma$ in rad/s/T, the magnetization energy
per cell is $E_{\text{cell}} = -\Ms \Vcell \, \mb{m} \cdot \Heff$ so
that the thermal-equilibrium probability density is
$\propto \exp(-E_{\text{cell}}/\kB T)$. The Fokker--Planck operator
associated with~\eqref{eq:sllg} has stationary solution equal to this
Boltzmann distribution if and only if the noise variance
$\sigma^2$ obeys the Einstein--type relation
%
\begin{equation}
\sigma^2 = \frac{2\, \alpha\, \kB T}{\gamma\, \Ms\, \Vcell},
\end{equation}
%
identical to~\eqref{eq:sigma}. This is the standard
fluctuation--dissipation result for Brown's stochastic LLG
equation~\cite{brown1963,garcia1998}. The factor of $\Vcell = a^2
\tCo$ converts the macrospin formula to a per-cell amplitude
appropriate for the discretized lattice; this is why a finer lattice
at the same $T$ gives a smaller per-cell field but the same
integrated-energy fluctuations.

%======================================================================
\begin{thebibliography}{9}

\bibitem{brown1963}
W.~F.~Brown, ``Thermal fluctuations of a single-domain particle,''
\emph{Physical Review} \textbf{130}, 1677 (1963).

\bibitem{garcia1998}
J.~L.~Garc\'ia-Palacios and F.~J.~L\'azaro,
``Langevin-dynamics study of the dynamical properties of small
magnetic particles,''
\emph{Physical Review B} \textbf{58}, 14937 (1998).

\bibitem{pham2024}
V.~T.~Pham et al., ``Fast current-induced skyrmion motion in
synthetic antiferromagnets,'' \emph{Science} (2024).

\bibitem{litzius2017}
K.~Litzius et al., ``Skyrmion Hall effect revealed by direct
time-resolved X-ray microscopy,'' \emph{Nature Physics} \textbf{13},
170 (2017).

\bibitem{rohart2016}
S.~Rohart, J.~Miltat, A.~Thiaville,
``Path to collapse for an isolated N\'eel skyrmion,''
\emph{Physical Review B} \textbf{93}, 214412 (2016).

\bibitem{kidger2021}
P.~Kidger, ``On Neural Differential Equations,'' PhD thesis,
University of Oxford, 2021 (\emph{diffrax} library).

\end{thebibliography}

\end{document}
